\newpage
\section{Reachability model checking in 2D space}
\label{ch:global}

% \section{Evaluation}

In this section we will analyze  another form of reachability that was described in 
Chapter \ref{ch:ModelCheckingLocal}. 
We will explore the reachability of various points on the field, assuming 
that the robot moves in steps, fixing the end effectors on the field one 
at a time.  

In our experimental system, the robot moves on 
a grid of discrete cells, picking up components and assembling a new 
robot. Using model checking we want to ensure that the robot can 
successfully gather the necessary parts to replicate, or otherwise detect 
scenarios leading to failure.

Formal methods can verify reachability of 
the goal state -- a completed new robot or a failure state, where replication 
becomes impossible and thereby increase confidence in the design. 
By analyzing counterexamples to these LTL properties, 
we can extract sequences of actions leading to success or failure.

\subsection{Model structure for 2D space reachability}

\subsubsection{Assumptions and simplifications}

To make the global reachability checking tractable, we model a simplified 
Self-Replicating Robotic System with several key assumptions and abstractions. 
These assumptions define the environment and the robot's capabilities in our NuSMV model:

\begin{enumerate}
    
    \item The robot operates on a two-dimensional grid of cells --a finite field.
        We consider only a single end-effector that occupies one cell on this grid at any time, 
    representing the robot's current position (with coordinates x 
    and y). The rest of the robot's body is not explicitly modeled on the 
    grid -- instead, the structure of the robot is represented abstractly by 
    a sequence of required part types. The end-effector 
    is initially anchored at an assembly site, taken to be the origin cell 
    (0,0) for convenience. All other objects (parts or obstacles) occupy discrete 
    cell positions on the grid. An example of a grid configuration is shown 
    on the diagram in Figure \ref{fig:AssembleDiagram}, 
    different part types are scattered on a field and the 
    robot must collect them.

    \begin{figure}[h!]
        \centering
        \includegraphics[width=12cm]{AssembleDiagram.png} 
        \caption{Simplified representation of a grid, robot and material parts.}
        \label{fig:AssembleDiagram}
    \end{figure}

\item The robot can move one cell at a time to a 
    neighboring cell (up, down, left, or right), and only into cells that 
    are unoccupied (have a type "None"). In other words, the robot's end-effector cannot 
    move into a cell that already contains a part or an obstacle. This models 
    a scenario where the robot must navigate around obstacles and cannot pass 
    through or occupy the same space as existing objects. Movement is 
    constrained to the four orthogonal directions on the grid (no diagonal moves). 
    Each move is treated as a discrete time step in the model.

\item  At any given position, the robot can "sense" or 
    inspect only its immediate neighbors for parts. The reach range for 
    checking is limited to the four adjacent cells. If a neighboring cell 
    contains a part, the robot can decide to check that part’s type. 
    This locality assumption simplifies sensing—there is no global visibility 
    of all parts at once, only local perception of adjacent cells.

\item When the robot finds a required part needed 
    for replication, it will pick up that part and carry it back to the 
    assembly location (the origin) to attach it to the new robot under construction.
    We abstract this entire retrieval and assembly process as a single action 
    that does not explicitly appear as a series of steps in the NuSMV model.
    In other words, once a part is found, the model assumes the robot instantly 
    returns to the assembly site with the part and the assembly index advances.
    This means the time and path taken to carry the part back are not 
    modeled; the effect (having assembled one more part) is reflected in 
    the state. This assumption significantly reduces model complexity by 
    not requiring us to simulate back-and-forth transport or the physical attachment 
    process. After assembly, the robot’s end-effector is conceptually back 
    at the assembly site (0,0) and ready to search for the next part. 
    The assembly site itself is not explicitly marked on the grid (we 
    assume it is at the starting position). We also assume that assembling the 
    part does not change the grid configuration in the model (the part is 
    not removed from the grid representation when taken, since the assembly 
    action is outside the modeled steps). 
\item In a real scenario, the part would be removed from the environment once 
    collected. However, by not altering the grid, we avoid introducing additional 
    transitions; we rely on the assembly index to track progress rather than 
    the physical absence of the part.
\end{enumerate}

Under these assumptions, the self-replication process can be seen as 
a sequential loop of searching for the next required part, checking if 
an encountered part matches the needed type, assembling the part if 
it matches, which implicitly resets the search to the next part, and 
continuing the search for subsequent parts until either the new robot 
is fully assembled or no further progress is possible. 

The high-level behavior can be represented by a state chart, 
as depicted in Figure \ref{fig:globalReachabilityStateChart}. 

\begin{figure}[h!]
    \centering
    \includegraphics[width=15cm]{GlobalReachabilityStatechart.png} 
    \caption{State chart of the self-replication system modeled for global reachability}
    \label{fig:globalReachabilityStateChart}
\end{figure}

In this state machine, the robot has distinct modes: MOVING, CHECKING,
ASSEMBLING, and two terminal states: FINISHED and FAILED. 

Initially, the robot is in the MOVING state, roaming the grid to find parts. 
When a neighbor cell is occupied by some object, the robot transitions to 
CHECKING and inspects that object’s type. 

If the object is of the type that the robot currently needs 
the next part in the assembly sequence, the robot enters the ASSEMBLING state,
meaning it picks up the part and 
(in one abstract step) delivers it to the assembly point and attaches it.

This successful assembly of a part causes the robot to increment an 
internal counter of how many parts of the sequence have been assembled. 

The robot then returns to MOVING to search for the next part, unless the 
sequence is now complete, in which case it transitions to FINISHED 
and means a new robot has been successfully constructed. 

If the robot is in CHECKING state and finds that the adjacent part is not the type 
it needs, it will go back to MOVING state and continue searching 
(the unwanted part is left in place). 

The process repeats until either the FINISHED state is reached 
(all required parts found in order) or a FAILURE condition occurs. 
A failure state is reached if the robot cannot 
find any required part and cannot move to any new location -- essentially, 
the robot is stuck with remaining parts needed, the assembly sequence isn't complete,
but with no available moves or accessible needed parts. 
In our model, this condition is captured by the situation where the 
robot is in MOVING state, still has parts to assemble, and every adjacent 
cell is either occupied (not NONE) by something (which might not be 
the needed part) or is out of bounds of the grid. In such a case, 
the model transitions the robot to FAILED, a terminal state indicating assembly 
cannot be completed. 

Both FINISHED and FAILED are terminal states in 
the model, that means once reached, the robot remains in that state. 
This state logic forms the basis of our NuSMV model.

\subsection{NuSMV Formal Model Structure}
\label{ch:globalNusmvStruct}
We constructed a formal model in NuSMV to capture the 
described behavior. The model is centered around the robot's 
position and state, the grid environment, and the assembly progress. 
In following sections we describe the main components of the model.

\subsubsection{State variables and initial conditions}
The NuSMV model defines discrete state variables corresponding to the 
robot and environment. Key variables include: grid, x and y coordinates,
robot\_state,sequence and checked\_cells.

     Grid or field is a 2D array representing the environment, 
        indexed from 0 to field\_max\_index in 
        both dimensions. Each cell in the grid can hold either “NONE” (if the 
        cell is empty) or a part type or obstacle type. We define a set of 
        possible cell values, e.g., NONE, TYPE1, TYPE2, TYPE3, OBSTACLE1,... . 
        In a given scenario, some of these types will be present and others 
        not, but the model is capable of representing any configuration of 
        parts. We initialize grid with a specific configuration where a handful 
        of cells have parts or obstacles and the rest are NONE. 
        An example of a small 4 x 4 size grid initialized as it shown 
        in Listing \ref{lst:globalReachGrid}:

    \begin{lstlisting}[language=NuSMV, caption={Initial representation of grid in NuSMV model}, 
label={lst:globalReachGrid}]
    initial_grid := [
    [NONE,	NONE,	NONE,	NONE,	NONE,	NONE],		
    [NONE,	NONE,	NONE,	NONE,	NONE,	NONE],		
    [NONE,	NONE,	NONE,	NONE,	NONE,	NONE],		
    [NONE,	NONE,	NONE,	TYPE2,	NONE,	NONE],		
    [TYPE1,	NONE,	NONE,	NONE,	NONE,	NONE],		
    [NONE,	NONE,	NONE,	NONE,	TYPE3,	NONE]
];
    \end{lstlisting}

    This example corresponds to a field where three part objects (
    of types TYPE1, TYPE2, TYPE3) are placed at specific coordinates, 
    and the rest are empty cells (NONE). An obstacle could similarly be 
    placed by setting a cell to OBSTACLE. 
    The robot's starting position is (0,0), which in this example is also the 
    assembly site and is empty. 
    
    2D coordinates x and y represent the robot's current position, integers ranging from 
    0 to field\_max\_index, where field\_max\_index = N-1 if the grid is N x N. 
    These coordinates represent the cell where the robot’s end-effector is 
    located.

    Robot state represented by robot\_state variable -- the mode of 
    operation of the robot, which can take values
    MOVING, CHECKING, ASSEMBLING, FINISHED, FAILED. This enumerated type 
    encodes the state machine described earlier.

    An integer variable assembly\_index tracking how many parts of the target 
    assembly sequence have been completed, or equivalently, the index of 
    the next required part in the sequence. 
    Initially, assembly\_index = 0 meaning no parts have 
    been assembled yet, and when assembly\_index exceeds robot\_size 
    the assembly is complete. In practice, reaching the last part triggers 
    the FINISHED state as described below.

    An array sequence,  holds the required part type for each index of 
    the robot's assembly sequence. This array is of length (robot\_size + 1). 
    For example, if the robot must assemble parts of types TYPE1, TYPE2,
    TYPE2, TYPE3 in that order, then robot\_size = 3 and sequence[0]=
    TYPE1, sequence[1]=TYPE2, sequence[2]=TYPE2, sequence[3]=TYPE3. 
    In the model initialization, we explicitly set each element of sequence to 
    define the blueprint of the robot we want to assemble. This sequence is 
    a parameter to the model, different robot designs can be specified by 
    changing these initial values.

     An auxiliary 2D array  checked\_cells (has the same dimensions as grid) used 
    to record which cells have been visited during the search. Each element 
    of checked\_cells can be either UNKNOWN or VISITED. Initially, all cells 
    are UNKNOWN. The intention is to mark cells as VISITED once 
    the robot has explored them, to prevent the robot from infinitely revisiting 
    the same empty cells in a loop. In the model, we will constrain movements 
    such that the robot does not move into a cell that is already marked VISITED.
    This prevents trivial cycles and helps the model checker focus on 
    new areas. The checked\_cells array is updated when the robot moves: when 
    transitioning to a new position, that position can be marked as visited.
    In our NuSMV model, this is encoded implicitly through transition constraints 
    rather than an explicit assignment.

At the beginning, we initialize the model variables as 
follows. The robot starts in state MOVING, with x = 0 and y = 0 (
assembly point). The assembly index is 0, which means no parts assembled yet. 
The grid is set to the predefined configuration of parts and obstacles for 
that scenario. 
The sequence array is initialized with the desired part 
sequence for the robot to build. All cells in checked\_cells start as 
UNKNOWN, except the starting cell which can be considered implicitly 
visited.


\subsubsection{State transition}

The core of the model is the state transition logic, which we encode using 
ASSIGN and TRANS blocks. 
We break down the transitions by robot state, and the simplified transition system
shown in Figure \ref{fig:globalReachabilityTransitionSystem}

\begin{figure}[h!]
    \centering
    \includegraphics[width=8cm]{GlobalReachabilityTransitionSystem.png} 
    \caption{Simplified transition system relative to robot\_state variable}
    \label{fig:globalReachabilityTransitionSystem}
\end{figure}

In the NuSMV model, the logic with robot\_state  is realized with a structured 
case statement for next(robot\_state) in ASSIGN section.
 When the robot is in MOVING state, it is actively changing the 
coordinates. We implement two key behaviors for MOVING:

If the robot is in MOVING and a required part is adjacent (in any of 
the four neighbors), then it should transition to CHECKING, 
as it shown in Listing \ref{lst:globalMoveCheck}. 
    \begin{lstlisting}[language=NuSMV, caption={Asign for transition from MOVING to CHECKING}, 
label={lst:globalMoveCheck}]
  (robot_state = MOVING & 
   ((x > 0                & grid[x - 1][y] != NONE) | 
    (x < field_max_index  & grid[x + 1][y] != NONE) | 
    (y > 0                & grid[x][y - 1] != NONE) | 
    (y < field_max_index  & grid[x][y + 1] != NONE))
  ) : CHECKING;
    \end{lstlisting}

% If any of these is true, a part (or obstacle) 
% is immediately adjacent, triggering a state change to CHECKING. This 
% rule corresponds to the robot noticing something to inspect.

If the robot is in MOVING and no required part can ever be found,
the robot should give up and enter FAILED, this part shown in Listing \ref{lse:globalMoveFail}. 
The condition we use to approximate "no required part exists or no 
moves possible" is: the robot is MOVING, still needs parts (assembly\_index <= 
robot\_size), and all adjacent cells are blocked (no empty neighbor to 
move into). 

    \begin{lstlisting}[language=NuSMV, caption={Assign for transition MOVING to FAILED}, 
label={lse:globalMoveFail}]
  (robot_state = MOVING & 
  (assembly_index <= robot_size) & 
  ((x = 0                 | grid[x - 1][y] != NONE ) & 
    (x = field_max_index  | grid[x + 1][y] != NONE) & 
    (y = 0                | grid[x][y - 1] != NONE ) & 
    (y = field_max_index  | grid[x][y + 1] != NONE ))
  ) : FAILED;
    \end{lstlisting}

If at the current position none of the four directions offers an empty 
cell (i.e. the robot is surrounded by occupied cells or edges), then 
the robot cannot move. 
captures dead-end situations (e.g., the robot is boxed in by obstacles or 
irrelevant parts that are not the one it needs).

The actual movement to a new cell is 
not decided in the next (robot\_state), instead, movement is governed by
separate TRANS constraints described later. 
In essence, if the robot stays in MOVING for the next step, we allow 
its (x, y) to change to an adjacent empty cell that hasn’t been visited 
yet. We encoded movement using additional transition constraints rather 
than in the main case, to separate the logic of where it moves from when 
it moves in the state machine.

When in robot\_state is in CHECKING, the robot is examining a specific 
neighboring part's type to see if it's what we need, the formal description is in Listing \ref{lst:globalCheck}. 

    \begin{lstlisting}[language=NuSMV, caption={Assign for transition from CHECKING state and changing of 
        auxiliary partX and partY variables}, 
    label={lst:globalCheck}]
    (robot_state = CHECKING & 
    assembly_index <= robot_size) : case
    grid[partX][partY]  = sequence[assembly_index] : ASSEMBLING; 
    grid[partX][partY] != sequence[assembly_index] : MOVING;
    esac;
    \end{lstlisting}
We need to decide whether to go to ASSEMBLING ,if the part matches the requirement
or back to MOVING if it does not.

To formalize this, we introduced variables 
partX and partY to record the coordinates of the part being checked, the conditions are
in Listing \ref{lst:globalPartXY}. 

    \begin{lstlisting}[language=NuSMV, caption={Assign for transition from CHECKING state and changing of 
        auxiliary partX and partY variables}, 
    label={lst:globalPartXY}]
    next(partX) := case
      ( robot_state = MOVING & next(robot_state) = CHECKING & 
        x > 0 & grid[x - 1][y] = sequence[assembly_index]
      ) : x - 1;  -- found in left neighbour cell (
        robot_state = MOVING & next(robot_state) = CHECKING & 
        x < field_max_index & grid[x + 1][y] = sequence[assembly_index]
      ) : x + 1;  -- found in right neighbour cell
      TRUE : x;  -- otherwise, retain previous (or irrelevant outside checking context)
    esac;
    next(partY) := case
      ( robot_state = MOVING & next(robot_state) = CHECKING & 
        y > 0 & grid[x][y - 1] = sequence[assembly_index]
      ) : y - 1;  -- found in upper neighbour cell
      ( robot_state = MOVING & next(robot_state) = CHECKING & 
        y < field_max_index & grid[x][y + 1] = sequence[assembly_index]
      ) : y + 1;  -- found in down neighbour cell
      TRUE : y;
    esac;
    \end{lstlisting}

When the robot transitions from MOVING to CHECKING, we simultaneously set 
partX and partY to the coordinates of the neighbor that it decided to 
inspect. In the model, we achieve this by writing next(partX) and next
(partY) with conditions mirroring the state change. 

For example, one rule is: if robot\_state = MOVING and next(robot\_state) = CHECKING and 
the left neighbor cell (x-1, y) contains the needed part type (grid[
x-1][y] = sequence[assembly\_index]), then next(partX) = x-1 and next
(partY) = y. Similar rules set partX and partY for the right, up, 
or down neighbor if that neighbor is the target part. 

If the part at (partX, partY) matches the required type for the current 
assembly index (i.e., grid[partX][partY] = sequence[assembly\_index]), 
then the robot will proceed to pick it up and assemble it.  

If the part is not the needed type (grid[partX][partY] != sequence[
assembly\_index]), then the robot will return to MOVING to keep searching. 
So we have next(robot\_state) = MOVING for that case. We coded this using 
a nested case within the CHECKING branch in NuSMV (so that from CHECKING
, the next state is a sub-case that depends on a comparison of the part
’s type).

The ASSEMBLING state represents that the robot 
has successfully fetched a part and is adding it to the assembly. 
We consider that the assembly of one part is completed within this state,
so the next step depends on whether there are more parts to assemble. 
The implementation of this transition shown in the Listing \ref{lst:globalAssembling}

\begin{lstlisting}[language=NuSMV, caption={Assign for transitions from ASSEMBLING state}, 
    label={lst:globalAssembling}]
    (robot_state = ASSEMBLING
    ) : case 
    assembly_index < robot_size : MOVING;
    assembly_index = robot_size : FINISHED;
    esac;
\end{lstlisting}

If the robot was assembling a part, and it was not the final part (
i.e. assembly\_index < robot\_size before assembly), then after assembling, 
the robot should go back to MOVING to search for the next part. We 
model this by next(robot\_state) = MOVING when robot\_state = ASSEMBLING and 
assembly\_index < robot\_size. Additionally, we will update the assembly\_index: while 
in ASSEMBLING, if it’s not the last part, we increment assembly\_index by 
1 for the next state (so the needed part type will advance to the next 
in the sequence).

If the robot was assembling and this part was the final one needed (
assembly\_index = robot\_size, meaning we just placed the last part), 
then the assembly is now complete. 
Therefore, we transition to FINISHED.
In this scenario, we also increment assembly\_index (
although it’s not strictly necessary once we mark FINISHED, but it 
could be used to indicate assembly\_index = robot\_size+1 meaning “
done”).

In summary, ASSEMBLING is a transient state: it will always go either to 
MOVING (with the next part index) or to FINISHED, in one step.

FINISHED and FAILED are the absorbing terminal states. We encode in Listing \ref{lst:globalFinFail} 
that if the robot is already in FINISHED, it stays in FINISHED, and similarly for FAILED. 

\begin{lstlisting}[language=NuSMV, caption={Terminal conditions for FINISHED and FAILED}, 
    label={lst:globalFinFail}]
      robot_state = FINISHED  : FINISHED;
      robot_state = FAILED    : FAILED;
\end{lstlisting}

Thus, once a terminal condition is reached, no further transitions occur 
and the trace effectively ends.



In addition to the robot\_state updates, the model includes constraints on 
how the x and y coordinates can change between steps. 
The movement and visit tracking logic described in TRANS block of the logic. 

If the robot will remain (or go) in the MOVING state at the next step
then we require that next(x), next(y) is one of the neighboring cells 
of (x, y) that is empty and not yet visited, this part shown  in Listing \ref{lst:globalXY}. 

\begin{lstlisting}[language=NuSMV, caption={Transition of x and y coordinates during MOVING}, 
    label={lst:globalXY}]
  TRANS 
    (next(robot_state) = MOVING ) -> (   
        ( next(x) = x - 1 & x > 0  
        & next(y) = y 
        & grid[x - 1][y] = NONE
        & checked_cells[x - 1][y] != VISITED
        ) |
        ( next(x) = x + 1 & x < field_max_index  
        & next(y) = y 
        & grid[x + 1][y] = NONE
        & checked_cells[x + 1][y] != VISITED 
        ) |
        (next(y) = y - 1 & y > 0 
        & next(x) = x 
        & grid[x][y - 1] = NONE
        & checked_cells[x][y - 1] != VISITED
        ) |
        (
        next(y) = y + 1 & y < field_max_index  
        & next(x) = x 
        & grid[x][y + 1] = NONE
        & checked_cells[x][y + 1] != VISITED 
        ));
\end{lstlisting}

 This big disjunction covers each of the four directions. 
 It essentially means: if we are continuing to move, the robot must move 
 by one step into a free cell that has not been visited before. 
It prevents illegal moves (e.g., jumping multiple cells or moving into occupied cells) and also 
 prevents staying in place during a MOVING step.

If the robot's next state is not MOVING (next(robot\_state) != MOVING), 
that implies the robot is engaging in checking or assembling, so we 
constrain that the robot’s coordinates do not change in that step, see Listing \ref{lst:globalMoveXY}

\begin{lstlisting}[language=NuSMV, caption={Condition that robot position doesnt change if not MOVING}, 
    label={lst:globalMoveXY}]
    TRANS (next(robot\_state) != MOVING) -> (next(x) = x & next(y) = y)
\end{lstlisting}

This ensures that during a CHECKING or ASSEMBLING step, the robot stays at the same 
location.

We also add TRANS (next(grid) = grid) because we assume the environment's
configuration of objects doesn't spontaneously change between steps 
(no external agent moves the parts or obstacles during the process). 
The only change to the grid would theoretically 
be removal of a part that was picked up, but as per our assembly abstraction,
we are not updating the grid for that, so indeed the grid remains static 
in the model. Similarly, TRANS (next(sequence) = sequence) is added because 
the required sequence doesn’t change over time (it's a fixed goal).


For completeness, we enforce that unless in ASSEMBLING (where assembly\_index increments),
the assembly\_index stays the same across a step: 
TRANS (robot\_state != ASSEMBLING) -> (next(assembly\_index) = assembly\_index). 
Combined with the earlier assignment rule that in ASSEMBLING state we do increment, 
this ensures assembly\_index only goes up exactly when a part is being 
assembled.

In our implementation, we simplified matters by 
only constraining moves to not go into VISITED cells, and we implicitly assume 
the robot doesn't immediately revisit its last position because it's now marked visited.
The initial cell (0,0) could be considered visited 
from the start to avoid the robot just stepping out and back in repeatedly.
In practice, because the robot always has the option to move into an 
unvisited cell (and the model checker tends to find a path to the goal 
if it exists, rather than looping), we found that the first counterexample trace 
it finds is typically a simple path without revisiting. 
Nonetheless, a more rigorous model would include an explicit update of checked\_cells[
next(x)][next(y)] := VISITED and possibly mark the start as visited in 
the init section for completeness.


\subsubsection{Specification for checking 2D reachability}

To perform reachability analysis of the terminal states (FINISHED or FAILED), 
we frame the problem as checking LTL properties that assert that those 
states never happen. By checking the negation of the reachability, 
the model checker will provide a counterexample path 
when the state is in fact reachable. 
We used two LTL formulas, first - successful assembly never occurs, in Listing \ref{lst:globalFin}

\begin{lstlisting}[language=NuSMV, caption={Terminal conditions for FINISHED and FAILED}, 
    label={lst:globalFin}]
    LTLSPEC G F (robot_state != FINISHED). 
\end{lstlisting}

This LTL formula reads as "always eventually, the robot\_state is not 
FINISHED," which is a somewhat indirect way of saying "robot\_state is 
never stuck in FINISHED forever." 

If this property is false, it means there exists a path where eventually robot
\_state = FINISHED and remains FINISHED thereafter (since once finished, 
it stays). In simpler terms, a counterexample to this property is 
a finite trace that ends in 
the FINISHED state (the model checker will actually produce a loop or 
repetition after FINISHED, but the interesting part is up to reaching FINISHED). 
Thus, a counterexample to this formula provides a concrete sequence of 
actions leading the robot to successfully assemble the new robot.

Another possible specification shown in Listing \ref{lst:globalFail}
\begin{lstlisting}[language=NuSMV, caption={Terminal conditions for FINISHED and FAILED}, 
    label={lst:globalFail}]
    LTLSPEC G F (robot\_state != FAILED). 
\end{lstlisting}
This formula asserts "it is always eventually not FAILED," i.e. the 
robot never gets permanently stuck in FAILED. If this is false, there 
is a scenario where the robot ends up in FAILED state. The counterexample trace 
for this formula will show how the robot ended up unable to complete the 
assembly.

In our NuSMV model file, we typically use first of these LTL specifications.
For example, to find a successful path in 2D.
But if we want to find a potential failure path, we check the second. 

In all interesting cases for our scenarios, the properties are 
false because we specifically set up scenarios where eventually the 
robot does finish (we ensure all needed parts are actually present somewhere 
on the grid) or could fail (if we design a scenario where a part is 
absent or unreachable). 

We leverage these traces as a source of information for 
analyzing and interpreting the robot behavior.
% , as described in the Chapter \ref{ch:validation}.


\subsection{Model generation and verification}

Writing a NuSMV model for one specific grid configuration and robot sequence is 
useful for a single-case analysis, but our goal was to explore many scenarios 
(different field sizes, different placements of parts, different robot 
part sequences) without manually coding each model variant. To facilitate this,
we developed a Python toolchain to automate model generation, execution,
and counterexample processing. 

The workflow can be summarized in the following steps:

We prepare a general NuSMV model template based on the  NuSMV model 
described in section \ref{ch:globalNusmvStruct}.
with placeholders for scenario-specific parameters. 

The template, stored in a file Assemble\_initial\_template.smv,
presented in Appendix \ref{ch:globalTemplate}, 
it includes formatting placeholders such 
as {field\_max\_index}, {object\_types}, {init\_robot\_sequence}, {
initial\_grid}, etc. in place of concrete values.
This template captures the structure of the model (variables, state transition 
logic, LTL specification), leaving certain details to be filled in.

We built a simple graphical user interface,shown in Figure \ref{fig:globalGUI1}, 
using Tkinter library, to allow interactive configuration: the user can specify the field size 
and robot length.
Then user have to press button "Generate field".

\begin{figure}[h!]
    \centering
    \includegraphics[width=9cm]{globalGUI1.png} 
    \caption{Graphical interface for 2D reachability checking}
    \label{fig:globalGUI1}
\end{figure}

The program opens the grid with buttons for each cell, where the user can select 
a cell's content from a dropdown (e.g. None, or specific part types), 
this part shown in Figure \ref{fig:globalGUI2}. 

The robot’s required part sequence can be set via dropdowns in a side 
panel. This interface makes it easy to set up custom scenarios without directly 
editing code.
The tool ensures the field and sequence lengths are consistent (it 
will create as many sequence slots as the robot size requires, etc.).

\begin{figure}[h!]
    \centering
    \includegraphics[width=15cm]{globalGUI2.png} 
    \caption{Graphical interface for 2D reachability checking with field}
    \label{fig:globalGUI2}
\end{figure}



Once the scenario is configured, user press "Generate model" button and
the  program uses the provided inputs to generate a concrete NuSMV model file.
This is done by taking the template and substituting in the actual values. 
The Python class takes user input or configuration for the scenario, 
such as the field size, the list of part types and obstacle types to 
consider, the initial positions of all objects on the field, and the 
robot's assembly sequence (which defines robot\_size and the sequence array 
values). 

The set of object types is derived from all items present on the 
grid (plus the “NONE” empty type); 
The robot's part sequence (a list like ["TYPE1", "TYPE2", "TYPE2", "TYPE3"]) 
is used both to set robot\_size (length minus 
one) and to fill in the init(sequence[i]) lines in the model. The initial 
grid configuration is formatted as a matrix literal for NuSMV. 

This model file is essentially the same structure for any scenario,
but with different initial conditions and parameters reflecting that 
scenario. The Python generator prints a summary to the user, confirming the 
model generation and showing the initial grid and sequence it encoded, 
for verification purposes.

After generating the model, with "Verify model" button we verify it using NuSmv. 
We automate this by invoking the nuSmv binary through a subprocess 
call on the generated model file.
The output will indicate whether the LTL specification was satisfied or 
violated. 

The significant part of our program is parsing the NuSmv counterexample to extract the 
robot's movement path and key events. 
We focus on the variables of interest: x, y, and the robot\_state (particularly 
noting when it reaches FINISHED or FAILED). 
Our Python script reads the 
output line by line and filters those lines that indicate a change or 
value of x or y or mention the robot\_state becoming FINISHED. 
By doing this, we reconstruct the sequence of coordinates the robot went through. 

The trace format usually lists assignments in 
a single step as separate lines.
By capturing x and y and assembling them in order, we build a list of 
(x, y) positions visited at each step of the counterexample path. 
The path indicates the cells the robot effector occupied 
over time.
The Python GUI then takes this path of coordinates and visually 
draws it on the grid for the user: connecting the sequence of cell centers 
with lines (in chronological order) so one can see the trajectory. 
The output text field of the GUI also shows a cleaned-up version of 
the nuSmv output and/or the path coordinates.

By automating this workflow, we can iterate quickly over different scenarios 
and immediately get the robot’s path if assembly is possible, or see 
where it fails. The counterexample trace effectively serves as a simulated 
execution of the self-replication process under the given scenario. 
Thus, we explore the state space with model checker  and provide an 
example execution leading to the outcome of interest.






